\section*{Bab \ref{ch:b1:first-law} --- Hukum Pertama Termodinamika}

\begin{solution}{exo:b1:first-law:1}
$Q = C_{P,m}\Delta T = 29.1 \times 100 = \qty{2.91}{kJ}$; $W = -nR\Delta T =
\qty{-0.83}{kJ}$; $\Delta U = Q + W = \qty{2.08}{kJ} = C_{V,m}\Delta T$.
\end{solution}

\begin{solution}{exo:b1:first-law:2}
$W = -nRT\ln(V_2/V_1) = -2494\ln(0.25) = \qty{+3.46}{kJ}$; $\Delta U = 0$ sehingga
$Q = \qty{-3.46}{kJ}$ (diberikan kepada rendamannya).
\end{solution}

\begin{solution}{exo:b1:first-law:3}
$T_2 = T_1 \times 10^{\gamma - 1} = 300 \times 10^{0.4} = \qty{754}{K}$; $P_2 =
10^{1.4} = \qty{25}{bar}$.
\end{solution}

\begin{solution}{exo:b1:first-law:4}
Argon: $C_{V,m} = \tfrac32 R = 12.5$, $C_{P,m} = \tfrac52 R = 20.8$ J/(\omterm{def:b1:kinetic-theory:scales}{mol} K),
$\gamma = 1.67$; nitrogen: $20.8$, $29.1$, $\gamma = 1.40$. Bagi air,
pemuaiannya ketika dipanaskan amat kecil, sehingga usaha $P\,\dd V$ pada tekanan tetap
dapat diabaikan terhadap \omterm{def:b1:first-law:heat}{kalornya}: $c_P - c_V \approx 0$.
\end{solution}

\begin{solution}{exo:b1:first-law:5}
$T_f = (200 \times 80 + 300 \times 20)/500 = \qty{44}{\degreeCelsius}$. Dengan
cangkirnya: $(200 \times 4.18 \times 80 + (300 \times 4.18 + 80) \times 20)/(500 \times
4.18 + 80) = \qty{43}{\degreeCelsius}$.
\end{solution}

\begin{solution}{exo:b1:first-law:6}
Esnya: hangatkan sampai $0$: $0.050 \times 2100 \times 10 = \qty{1.05}{kJ}$; lalu cairkan:
\qty{16.7}{kJ}; totalnya \qty{17.8}{kJ}. Airnya dapat memberikan \qty{25.1}{kJ} sampai
$0$: seluruhnya mencair; lalu $0.250 \times 4180\,T_f = 25.1 - 17.8$ kJ: $T_f =
\qty{7}{\degreeCelsius}$. Dengan \qty{150}{g}: dituntut \qty{53.3}{kJ} $>$ \qty{25.1}{kJ}:
hanya $(25.1 - 3.15)/334 = \qty{66}{g}$ es yang mencair, $T_f = \qty{0}{\degreeCelsius}$,
\qty{84}{g} es tersisa.
\end{solution}

\begin{solution}{exo:b1:first-law:7}
$W = -P_1(V_2 - V_1) + 0 - P_2(V_1 - V_2) + 0 = (P_2 - P_1)(V_2 - V_1) > 0$:
siklusnya ditempuh berlawanan arah jarum jam (memuai pada tekanan rendah,
memampat pada tekanan tinggi), sehingga gasnya \emph{menerima} usaha total, sebesar
luas yang dilingkupinya; searah jarum jam ia akan menyampaikannya --- sebuah mesin.
\end{solution}

\begin{solution}{exo:b1:first-law:8}
$T_2 = 300 \times 7^{0.286} = \qty{523}{K}$. Udara mampat yang panas memanaskan
tabungnya secara konduksi pada tiap langkah. Di dalam bannya udaranya mendingin ke
suhu sekitarnya pada volume tetap: $P \propto T$, tekanannya turun menurut nisbah
suhunya --- tambahkanlah lagi.
\end{solution}

\begin{solution}{exo:b1:first-law:9}
$W = 0$ (tak ada tekanan luar untuk dilawan), $Q = 0$: $\Delta U = 0$,
sehingga $\Delta T = 0$ bagi \omterm{def:b1:kinetic-theory:model}{gas ideal}. Tak reversibel: gasnya tak pernah mengalir
kembali dengan sendirinya. Gas nyata, yang molekulnya saling menarik, mendingin sedikit
(sebagian \omterm{thm:b1:work-and-energy:ke}{energi kinetiknya} menjadi \omterm{def:b1:work-and-energy:conservative}{energi potensial} ketika mereka saling menjauh).
\end{solution}

\begin{solution}{exo:b1:first-law:10}
$Q = L_v = \qty{2.26}{MJ}$; $W = -P(V_g - V_l) = -1.013\times10^5 \times 1.67 =
\qty{-0.17}{MJ}$; $\Delta U = \qty{2.09}{MJ}$: $93\%$ \omterm{def:b1:first-law:heat}{kalornya} dipakai untuk
memisahkan molekulnya (\omterm{def:b1:kinetic-theory:U}{energi dalam}), $7\%$ untuk mendorong mundur
atmosfernya.
\end{solution}

\begin{solution}{exo:b1:first-law:11}
$W = -P_2(V_2 - V_1)$, $\Delta U = nC_{V,m}(T_2 - T_1)$, $V_2 = nRT_2/P_2$, $V_1
= nRT_1/P_1$: $\tfrac52(T_2 - 300) = -T_2 + 300 \times 5$, $3.5T_2 = 2250$,
$T_2 = \qty{643}{K}$, $V_2 = \qty{10.7}{L}$. Reversibel: $T_2 = 300 \times
5^{0.286} = \qty{475}{K}$, $V_2 = \qty{7.9}{L}$. Pemampatan yang tiba-tiba melakukan
lebih banyak usaha pada gasnya (seluruh \qty{5}{bar} bekerja sejak awal) dan berakhir
lebih panas serta kurang termampatkan.
\end{solution}

\begin{solution}{exo:b1:first-law:12}
$\rho = PM/RT = \qty{1.20}{kg/m^3}$: $\sqrt{\gamma P/\rho} = \sqrt{1.4 \times
1.013\times10^5/1.20} = \qty{343}{m/s}$; $\sqrt{P/\rho} = \qty{290}{m/s}$, $15\%$
terlalu kecil. Dengan $P/\rho = RT/M$: $c = \sqrt{\gamma RT/M}$; berbanding $u = \sqrt{3RT/M}$
= \qty{502}{m/s}: $c/u = \sqrt{\gamma/3} = 0.68$.
\end{solution}

\begin{solution}{pb:b1:first-law:1}
\textbf{1.} $n = P_1V_1/RT_1 = 10^5 \times 2.12\times10^{-4}/2494 = \qty{8.5e-3}{mol}$.

\textbf{2.} $PV^\gamma$ tetap: $V_2 = V_1(P_1/P_2)^{1/\gamma} = 212 \times
7^{-0.714} = \qty{53}{cm^3}$.

\textbf{3.} $T_2 = T_1(P_2/P_1)^{(\gamma - 1)/\gamma} = 300 \times 7^{0.286} =
\qty{523}{K}$.

\textbf{4.} $W = nC_{V,m}(T_2 - T_1) = 8.5\times10^{-3} \times 20.8 \times 223 =
\qty{39}{J}$.

\textbf{5.} Udara yang harus ditambahkan: $\Delta n = \Delta P\,V/RT = 6\times10^5 \times 2\times10^{-3}/
2494 = \qty{0.48}{mol}$: $0.48/8.5\times10^{-3} = 57$ langkah, kira-kira
\qty{2.2}{kJ} --- semenit usaha yang sungguh-sungguh.

\textbf{6.} Pada volume tetap $P \propto T$: tekanannya turun ketika udaranya
mendingin (sampai sebesar $300/523$ bagi udara langkah terakhirnya); pesepedanya menambah
beberapa langkah setelah jeda.

\textbf{7.} $W_{\text{iso}} = nRT\ln(P_2/P_1) = 8.5\times10^{-3} \times 2494 \times
1.95 = \qty{41}{J} > \qty{39}{J}$: \omterm{def:b1:point-kinematics:frame}{lintasan} isotermalnya memampatkan gasnya
lebih jauh (ke \qty{30}{cm^3} dan bukan \qty{53}{cm^3}) untuk mencapai \qty{7}{bar},
karena ia tetap dingin --- volume yang disapu lebih besar, usahanya lebih besar.

\textbf{8.} $n = 10^5 \times 5\times10^{-4}/2494 = \qty{0.020}{mol}$; $T_2 =
300 \times 20^{0.4} = \qty{994}{K}$.

\textbf{9.} $P_2 = 20^{1.4} = \qty{66}{bar}$.

\textbf{10.} \qty{994}{K} jauh di atas titik nyala sendiri bahan bakarnya:
bahan bakar yang disuntikkan menyala begitu bersentuhan. Mesin bensin memampatkan
campuran udara--bahan bakarnya hanya sampai kira-kira \qty{750}{K} agar ia tak menyala
sebelum businya memercik; nisbah yang terlalu tinggi menimbulkan ``ketukan''.

\textbf{11.} $W = nC_{V,m}(T_2 - T_1) = 0.020 \times 20.8 \times 694 = \qty{289}{J}$.

\textbf{12.} $25 \times 289 = \qty{7.2}{kW}$ tiap tabung, yang tersimpan dalam gas
panasnya dan dikembalikan pada ekspansinya.

\textbf{13.} Ketakreversibelan (gasnya tak menuruti \omterm{def:b1:point-kinematics:frame}{lintasan}
kuasistatiknya; ada usaha tambahan yang terlesapkan) menaikkan $T_2$; \omterm{def:b1:first-law:heat}{kalor} yang hilang ke dindingnya
menurunkannya. Pada mesin yang sebenarnya biasanya yang kedua menang: $T_2$ sedikit
di bawah nilai idealnya.

\textbf{14.} $Q = 20\times10^{-6} \times 43\times10^6 = \qty{860}{J}$.

\textbf{15.} Isobarik: $Q = nC_{P,m}(T_3 - T_2)$: $T_3 - T_2 = 860/(0.020 \times
29.1) = \qty{1480}{K}$, $T_3 = \qty{2470}{K}$; $V_3 = V_2T_3/T_2 = 25 \times
2.48 = \qty{62}{cm^3}$.

\textbf{16.} $W = -P_2(V_3 - V_2) = -66\times10^5 \times 37\times10^{-6} =
\qty{-244}{J}$: usaha yang disampaikan gasnya.

\textbf{17.} $T_4 = T_3(V_3/V_1)^{\gamma - 1} = 2470 \times (0.124)^{0.4} =
\qty{1070}{K}$; $W = nC_{V,m}(T_4 - T_3) = 0.020 \times 20.8 \times (-1400) =
\qty{-582}{J}$.

\textbf{18.} Usaha total yang disampaikannya $= 244 + 582 - 289 = \qty{537}{J}$; nisbahnya
$537/860 = 62\%$ (efisiensi Diesel ideal pada setelan ini; mesin
yang sebenarnya mencapai kira-kira $40\%$).

\textbf{19.} $1 - 300/2470 = 88\%$: siklus Diesel ideal tinggal jauh
di bawah batas Carnot, karena \omterm{def:b1:first-law:heat}{kalornya} diterima sepanjang suatu rentang
suhu dan bukan pada suhu tertingginya.

\textbf{20.} Sepanjang satu siklus $\Delta U = 0$: $Q_{\text{keluar}} = -(860 - 537) =
\qty{-323}{J}$, yakni \qty{323}{J} pergi bersama gas buangnya; periksa:
$nC_{V,m}(T_4 - T_1) = 0.020 \times 20.8 \times 770 = \qty{320}{J}$.

\textbf{21.} Osilasinya cepat (sekitar sesekon) dibandingkan dengan konduksi
\omterm{def:b1:first-law:heat}{kalor} melalui \qty{10}{L} udara: adiabatik. Dengan melinearkan
$PV^\gamma$: $\dd P/P = -\gamma\,\dd V/V = -\gamma Sx/V$.

\textbf{22.} Gaya pada bolanya $S\,\dd P = -\gamma PS^2x/V$: $m\ddot x =
-(\gamma PS^2/V)x$, harmonik, $\omega_0^2 = \gamma PS^2/mV$.

\textbf{23.} $\omega_0^2 = 1.4 \times 10^5 \times 4\times10^{-8}/(0.0165 \times 0.010)
= \qty{33.9}{s^{-2}}$: $T = 2\pi/5.83 = \qty{1.08}{s}$. Dari $T = \qty{1.10}{s}$:
$\gamma = 4\pi^2mV/(T^2PS^2) = 4\pi^2 \times 0.0165 \times 0.010/(1.21 \times 10^5
\times 4\times10^{-8}) = 1.35$.

\textbf{24.} Gesekan meredam geraknya dan, bila ia kering, menggeser
kedudukan setimbangnya; kebocoran di sisi bolanya melepaskan udara selama pemampatan,
sehingga gaya pemulihnya mengecil: keduanya memperpanjang periodenya dan membuat
$\gamma$ berat sebelah ke bawah --- seperti $1.35$ di atas.

\textbf{25.} Pompa dan diesel: $Q = 0$, sehingga $\Delta U = W$ --- usaha menjadi
\omterm{def:b1:kinetic-theory:U}{energi dalam} dan suhu, menurut \omterm{prop:b1:first-law:transformations}{hukum Laplace} dengan eksponen
$\gamma$. Bolanya: pemampatan yang cepat bersifat adiabatik, dan $\gamma$ menetapkan
kekakuan sebuah pegas udara.
\end{solution}
